% BEGIN S1
\par\medskip\noindent\textbf{1 Introduction}\par
% BEGIN S1.p1
The Fréchet distance is a distance metric for polygonal curves, where each curve is represented as a piecewise-linear interpolation between its vertices. There are different variants of this metric and the two most common variants are the discrete and the continuous Fréchet distance. Alt and Godau [3] were the first to study the computational complexity of the Fréchet distance. Let $n$ and $m$ denote the number of vertices of the two curves, which we refer to as their complexities, with $n\geq m$. They showed that the continuous Fréchet distance can be computed in $\mathcal{O}(nm\log n)$ time. Eiter and Mannilla [21] defined the discrete Fréchet distance and gave an $\mathcal{O}(nm)$ time algorithm. Although both results were slightly improved later [15, 11, 1], Bringmann [6] showed that for $d$-dimensional curves with $d\geq 2$ it is very unlikely that significantly faster algorithms exists: unless the Strong Exponential Time Hypothesis (SETH) fails, no $(1.001)$-approximation algorithm for the continuous or discrete Fréchet distance can run in $\mathcal{O}((nm)^{1-\delta})$ time for any $\delta>0$. Later, Bringmann and Mulzer [10] extended this result for the discrete Fréchet distance in 1D and increased the approximation factor up to $1.4-\varepsilon$ for any $\varepsilon>0$111In [10], this result is stated only for $n=m$, but the same construction works for the case $n\neq m$ as well.. For the continuous and the discrete Fréchet distance in 1D, Buchin, Ophelders, and Speckmann [12] showed a hardness result for approximation factors up to $3-\varepsilon$, but only in the case where $n=m$. A recent result by Blank and Driemel [5] shows that the continuous Fréchet distance in 1D can indeed be computed faster when $m=n^{\alpha}$ with $\alpha\in(0,1)$, namely in $\mathcal{O}(n\log n+m^{2}\log^{2}n)$ time.
\par
\par
% END S1.p1
% BEGIN S1.SS0.SSS0.P0.SPx1
\par\medskip\noindent\textbf{Related Work.}\par
% BEGIN S1.SS0.SSS0.P0.SPx1.p1
In recent years, many works have studied approximation algorithms for the Fréchet distance [18, 26, 27]. This line of research culminated in the first constant-factor approximation algorithm for the Fréchet distance that runs in subquadratic time, due to Cheng, Huang, and Zhang [17]. Their randomized $(7+\varepsilon)$-approximation algorithm for the continuous Fréchet distance runs in $\mathcal{O}(nm^{0.99}\log(n/\varepsilon))$ time for any $\varepsilon>0$ and succeeds with probability at least $1-1/n^{6}$.
\par
\par
% END S1.SS0.SSS0.P0.SPx1.p1
% BEGIN S1.SS0.SSS0.P0.SPx1.p2
A simple idea to $3$-approximate the Fréchet distance between curves where $m=n^{\alpha}$ with $\alpha\in(0,1)$ in subquadratic time is the following. Compute a curve $P^{\prime}$ of complexity $n^{\prime}\in o(n)$ and its distance $\delta$ to $P$ such that every curve $Q$ of complexity $m$ has Fréchet distance at least $\delta$ to $P$. Then compute the Fréchet distance $\delta^{\prime}$ between $P^{\prime}$ and $Q$. Using triangle inequality, one can show that the Fréchet distance between $P$ and $Q$ is at least $\max\{\delta^{\prime}/2,\delta\}$ and at most $3\cdot\max\{\delta^{\prime}/2,\delta\}$, which yields a $3$-approximation. The difficulty of this approach is to compute such a curve $P^{\prime}$. Two related variants of this simplification problem have been studied. The first is computing an optimal $m$-simplification of $P$, that is, a curve $P^{*}$ of complexity at most $m$ minimizing the Fréchet distance to $P$ among all curves of complexity at most $m$. The second is computing an optimal $\delta$-simplification, that is, a curve $P^{*}$ of minimum complexity among all curves that have Fréchet distance at most $\delta$ to $P$. However, no previous result achieved a subquadratic running time for curves in arbitrary dimension when the approximation factor on the complexity of the simplified curve is constant [2, 4, 7, 13, 16, 22, 28].
\par
\par
% END S1.SS0.SSS0.P0.SPx1.p2
% BEGIN S1.SS0.SSS0.P0.SPx1.p3
A different approach to reduce the running time of computing the Fréchet distance is via Fréchet distance oracles. Such an oracle is a data structure that preprocesses an input curve and subsequently answers Fréchet distance queries with respect to query curves. Driemel and Har-Peled [19] described a data structure of size $\mathcal{O}(n\log n)$ using $\mathcal{O}(m^{2}\log n\log(m\log n))$ query time, which approximates the continuous Fréchet distance up to a (large) constant factor. For exact queries, [14] presented a data structure of size $\mathcal{O}(nm)^{\text{poly}(d,m)}$ and query time in $\mathcal{O}((md)^{\mathcal{O}(1)}\log(nm))$, where the query complexity $m$ is given at preprocessing time. Blank and Driemel [5] showed that for curves in 1D, there is a data structure of size $\mathcal{O}(n\log n)$ that computes the continuous Fréchet distance exactly in $\mathcal{O}(m^{2}\log^{2}n)$ time. For the discrete Fréchet distance, Filtser and Filtser [22] improved [20] and presented a data structure with size in $\mathcal{O}(1/\varepsilon)^{kd}\log(\varepsilon)^{-1}$ that approximates the discrete Fréchet distance within a factor of $(1+\varepsilon)$ using expected $\widetilde{\mathcal{O}}(kd)$ query time, where $\widetilde{\mathcal{O}}(\cdot)$ hides polylogarithmic factors in $n$. Gudmundsson, Seybold, and Wong [24] used the lower bound construction by Bringmann [6] to show in a similar fashion to [8, 20, 25] the following. In 2D, there exists no Fréchet distance oracle with $\textup{poly}(n)$ preprocessing time and $\mathcal{O}((nm)^{1-\delta})$ query time and approximation factor $1.001$ unless SETH fails.
\par
\par
% END S1.SS0.SSS0.P0.SPx1.p3
% BEGIN S1.T1
 Table 1: Upper and Lower Bounds for the Fréchet distance Computation. 
\noindent Upper Bounds\par
\noindent  | Discrete | Continuous\par
\noindent 1D | ${\mathcal{O}}(mn)$ | exact | [21, 1] | $\widetilde{\mathcal{O}}(mn)$ | exact | [3, 11, 15]\par
\noindent $\widetilde{\mathcal{O}}(n+m^{2})$ |  $2$-approx. | Cor.20  | $\widetilde{\mathcal{O}}(n+m^{2})$ | exact | [5]\par
\noindent dD | ${\mathcal{O}}(mn)$ | exact | [21, 1] | $\widetilde{\mathcal{O}}(mn)$ | exact | [3, 11, 15]\par
\noindent $\widetilde{\mathcal{O}}(n+m^{2})$ |  $(3+\varepsilon)$-approx. | Thm.23  | $\widetilde{\mathcal{O}}(n+m^{2})$ |  $(3+\varepsilon)$-approx. | Thm.23 \par
\noindent $\widetilde{\mathcal{O}}(nm^{0.99})$ |  $(7+\varepsilon)$-approx. | [17] | $\widetilde{\mathcal{O}}(nm^{0.99})$ |  $(7+\varepsilon)$-approx. | [17]\par
\noindent Lower Bounds\par
\noindent  | Discrete | Continuous\par
\noindent 1D | $\nexists\nobreak\ {\mathcal{O}}(m^{2-\delta})$ |  $(3-\varepsilon)$-approx. | [12] | $\nexists\nobreak\ {\mathcal{O}}(m^{2-\delta})$ |  $(3-\varepsilon)$-approx. | [12]\par
\noindent $\nexists\nobreak\ \widetilde{\mathcal{O}}((mn)^{1-\delta})$ |  $(1.4-\varepsilon)$-approx. | [10] |  |  | \par
\noindent $\nexists\nobreak\ \widetilde{\mathcal{O}}((mn)^{1-\delta})$ |  $(2-\varepsilon)$-approx. | Thm.9  |  |  | \par
\noindent 2D |  ${\nexists\nobreak\ \mathcal{O}}((mn)^{1-\delta})$ $(1.001)$-approx. [6] \par
\noindent  ${\nexists\nobreak\ \mathcal{O}}((mn)^{1-\delta})$ $(1+\sqrt{2}-\varepsilon)$-approx. in $L_{2}$ Thm.10 \par
\noindent  ${\nexists\nobreak\ \mathcal{O}}((mn)^{1-\delta})$ $(3-\varepsilon)$-approx. in $L_{\infty}$ Thm.10 \par
% END S1.T1
% END S1.SS0.SSS0.P0.SPx1
% BEGIN S1.SS0.SSS0.P0.SPx2
\par\medskip\noindent\textbf{Our Results.}\par
% BEGIN S1.SS0.SSS0.P0.SPx2.p1
Our first set of results shows lower bounds for algorithms that approximate the Fréchet distance conditioned on the Orthogonal Vectors Hypothesis (OVH). We show that the discrete Fréchet distance between 1-dimensional curves cannot be $(2-\varepsilon)$-approximated in $\mathcal{O}((nm)^{1-\delta})$ time for any $\varepsilon,\delta>0$ unless OVH fails. This increases the approximation factor of $1.4-\varepsilon$ by Bringmann and Mulzer [10]. Using a similar construction in 2D, we show that the continuous and discrete Fréchet distance cannot be approximated within a factor less than $1+\sqrt{2}$ (resp. $3$) when the curves lie in the Euclidean space (resp. in the $L_{\infty}$-space) in $\mathcal{O}((nm)^{1-\delta})$. This leads to lower bounds for Fréchet distance oracles as well and answers an open question by [24].
\par
\par
% END S1.SS0.SSS0.P0.SPx2.p1
% BEGIN S1.SS0.SSS0.P0.SPx2.p2
To obtain this first set of results, we reduce from the Orthogonal Vector Problem: Given two sets of vectors $U,V\subset\{0,1\}^{d}$ of size $m$ and $n$, do there exist two vectors $u\in U$ and $v\in V$ such that $\langle u,v\rangle=0$? Then, the results follow by OVH, which is implied by SETH [29, 9].
\par
\par
% END S1.SS0.SSS0.P0.SPx2.p2
% BEGIN Thmtheorem1
\par\medskip\noindent\textbf{Definition 1 (Orthogonal Vectors Hypothesis (OVH)).}\par
For all $\delta>0$ there exists a $c>0$ such that there is no algorithm solving OV instances $U,V\subset\{0,1\}^{d}$ of size $m$ and $n$ with $m\leq n$ and $d=c\log n$ in time $\mathcal{O}((nm)^{1-\delta})$.
\par
\par
\par
% END Thmtheorem1
% BEGIN S1.SS0.SSS0.P0.SPx2.p3
We construct a curve $P$ of complexity $\mathcal{O}(n+m)$ depending on $V$ and the size of $U$ and a curve $Q$ of complexity $\mathcal{O}(m)$ depending only on $U$. Then, we show that the Fréchet distance between $P$ and $Q$ is at most $1$ if and only if the OV instance has a solution. In order to get a lower bound also for the case where $n\gg m$, we construct $P$ such that the main part of $P$ lies in a ball of radius $1$. Hence, this main part can be matched to a single point of $Q$. Further, to obtain an inapproximability result of $2-\varepsilon$ in 1D, the points of $P$ and $Q$ have integer coordinates. Hence, there are only three different values that the vertices of the main part of $P$ can have. This significantly restricts our possibilities to define $P$ and is the main difference to the construction by Bringmann and Mulzer [10]. They also note in their paper that obtaining a lower bound for an approximation factor of $1.4$ or higher requires a significantly different construction compared to theirs.
\par
\par
% END S1.SS0.SSS0.P0.SPx2.p3
% BEGIN S1.SS0.SSS0.P0.SPx2.p4
On the positive side, we present a data structure that preprocesses a 1-dimensional curve of complexity $n$ in $\mathcal{O}(n\log n)$ time using $\mathcal{O}(m)$ space that $2$-approximates the discrete Fréchet distance to query curves of complexity $m$ in $\mathcal{O}(m^{2}\log m)$ time. Hence, there exists a $2$-approximation algorithm for the discrete Fréchet distance in $\mathcal{O}(n\log n+m^{2}\log m)$ time matching our approximation factor lower bound and essentially the time lower bound by [12]. Secondly, we match our approximation factor lower bound in the $L_{\infty}$-space by giving a simple $(3+\varepsilon)$-approximation algorithm for the discrete and continuous Fréchet distance using $\mathcal{O}((n+m^{2})\log n)$ time for curves in any dimension under any $L_{p}$-norm. Table 1 summarizes the results of this paper and related work.
\par
\par
% END S1.SS0.SSS0.P0.SPx2.p4
% END S1.SS0.SSS0.P0.SPx2
% END S1
% BEGIN S2
\par\medskip\noindent\textbf{2 Preliminaries}\par
% BEGIN S2.p1
For any two points $p_{1},p_{2}\in\mathbb{R}^{d}$, $\overline{p_{1}p_{2}}$ denotes the directed line segment connecting $p_{1}$ with $p_{2}$. A polygonal curve $P$ of complexity $n$ is formed by ordered line segments $\overline{P(i)P(i+1)}$ of points $P(1),P(2),\dotso,P(n)$, where $P(i)\in\mathbb{R}^{d}$. It can be viewed as a function $P:[1,n]\rightarrow\mathbb{R}^{d}$, where $P(i+\alpha)=(1-\alpha)P(i)+\alpha P(i+1)$ for $i\in\{1,\dotso,n\}$ and $\alpha\in[0,1]$. This curve is also denoted with $\langle P(1),\dotso,P(n)\rangle$. We call the points $P(i)$ vertices and the line segments $\overline{P(i)P(i+1)}$ edges of $P$. Further, $P[s,t]$ is the subcurve of $P$ obtained from restricting the domain to the interval $[s,t]$. Two curves $P=\langle P(1),\dots,P(n)\rangle$ and $Q=\langle Q(1),\dots,Q(m)\rangle$ can be composed into the curve $P\circ Q=\langle P(1),\dots,P(n),Q(1),\dots,Q(m)\rangle$. For a number $k\in\mathbb{N}_{0}$, define $P^{k}$ be the composition of $k$ copies of $P$.
\par
\par
% END S2.p1
% BEGIN S2.p2
Let $P:[1,n]\rightarrow\mathbb{R}^{d}$ and ${Q:[1,m]\rightarrow\mathbb{R}^{d}}$ be two polygonal curves. A continuous matching $M$ between $P$ and $Q$ is defined by two functions $h_{P}\in\mathcal{F}_{n}$ and $h_{Q}\in\mathcal{F}_{m}$ and consists of the tuples $(h_{P}(a),h_{Q}(a))$ for $a\in[0,1]$, where $\mathcal{F}_{n}$ is the set of all continuous, non-decreasing functions $h:[0,1]\rightarrow[1,n]$ with ${h(0)=1}$ and $h(1)=n$, respectively $\mathcal{F}_{m}$ for $m$. Then, the (continuous) Fréchet distance between $P$ and $Q$ is defined as
\par
\[\begin{aligned}
d_{\text{F}}(P,Q)=\min_{\text{continuous matching }M}\ \max_{(x,y)\in M}\|P(x)%
-Q(y)\|.
\end{aligned}\]
For a fixed $\delta$, the free space $F_{\delta}(P,Q)$ is a subset of $[1,n]\times[1,m]$ such that $(x,y)\in F_{\delta}(P,Q)$ if and only if $\|P(x)-Q(y)\|\leq\delta$. Points in $([1,n]\times[1,m])\setminus F_{\delta}(P,Q)$ are contained in the non-free space. Hence, the continuous Fréchet distance between $P$ and $Q$ is at most $\delta$ if a continuous matching $M\subset F_{\delta}(P,Q)$ between $P$ and $Q$ exists.
\par
\par
% END S2.p2
% BEGIN S2.p3
A discrete matching $M$ between $P$ and $Q$ is an ordered sequence that starts in $(1,1)$, ends in $(n,m)$ and for any consecutive tuples $(i,j)$, $(i^{\prime},j^{\prime})$ in $M$ it holds that $(i^{\prime},j^{\prime})\in\{(i+1,j),(i,j+1),(i+1,j+1)\}$. Then, the discrete Fréchet distance between $P$ and $Q$ is
\par
\[\begin{aligned}
d_{\text{dF}}(P,Q)=\min_{\text{discrete matching }M}\ \max_{(i,j)\in M}\|P(i)-%
Q(j)\|.
\end{aligned}\]
The free space matrix $M_{\delta}$ is an $n\times m$ matrix, where $M_{\delta}[i,j]=1$ if $|P(i)-Q(j)|\leq\delta$ and otherwise $M_{\delta}[i,j]=0$. We say that $(i,j)$ is reachable if $d_{\text{dF}}(P[1,i],Q[1,j])\leq\delta$.
\par
\par
% END S2.p3
% BEGIN Thmtheorem2
\par\medskip\noindent\textbf{Observation 2.}\par
Let $P_{1},P_{2},Q_{1},Q_{2}$ be curves. If $d_{\text{F}}(P_{1},Q_{1})\leq\delta$ and $d_{\text{F}}(P_{2},Q_{2})\leq\delta$, then ${d_{\text{F}}(P_{1}\circ P_{2},Q_{1}\circ Q_{2})\leq\delta}$. For the discrete Fréchet distance, the analogue statement holds.
\par
\par
\par
% END Thmtheorem2
% BEGIN Thmtheorem3
\par\medskip\noindent\textbf{Observation 3.}\par
Let $P$ and $Q$ be polygonal curves, then $d_{\text{F}}(P,Q)\leq d_{\text{dF}}(P,Q)$.
\par
\par
\par
% END Thmtheorem3
% END S2
% BEGIN S3
\par\medskip\noindent\textbf{3 Approximation lower bounds for the imbalanced case}\par
% BEGIN S3.p1
In this section, we show a lower bound for the discrete Fréchet distance in 1D and for the discrete and continuous Fréchet distance in 2D depending on OVH. The construction for both bounds is very similar and we start with giving a schematic view of the reduction.
\par
\par
% END S3.p1
% BEGIN S3.p2
Let $U,V\subset\{0,1\}^{d}$ be two sets of size $m$ and $n$, where $m\leq n$ and denote with $v_{ik}$ (resp. $u_{ik}$) the $k$-th bit of the $i$-th vector of $V$ (resp. $U$). We construct two 1-dimensional (resp. 2-dimensional) curves $P$ and $Q$ of complexity $\mathcal{O}(nd)$ and $\mathcal{O}(md)$, where $P$ depends on $V$ and $Q$ on $U$. We show that their discrete (and continuous) Fréchet distance is at most $1$ if and only if two vectors that are orthogonal exist in $V$ and $U$. Define
\par
\[\begin{aligned}
P=P_{s}\circ(P^{*})^{m-1}\circ P_{null}\circ\bigcirc_{i=1}^{n}\left((P_{0})^{d%
+1}\circ P_{1}\circ\bigcirc_{\ell=1}^{d}P_{v_{i\ell}}\circ P_{1}\right)\circ P%
_{null}\circ(P^{*})^{m-1}\circ P_{s},
\end{aligned}\]
\[\begin{aligned}
Q=(Q^{*})^{m-1}\circ\bigcirc_{j=1}^{m}\left(Q_{null}\circ Q_{c}\circ(Q_{1})^{d%
+1}\circ Q_{0}\circ\bigcirc_{\ell=1}^{d}Q_{u_{j\ell}}\circ Q_{0}\circ Q_{c}%
\right)\circ Q_{null}\circ(Q^{*})^{m-1},
\end{aligned}\]
where the subcurves are defined later. Since the complexity of $P$ can be arbitrarily larger than the complexity of $Q$, we construct the curve $P$ such that the main part of $P$, i.e., in between the two subcurves $P_{null}$, is contained in a ball of radius $1$. Then, we can match this whole part to a single point on $Q$. This point is $Q_{c}$ in our construction. At a high level idea, we want to enforce that in a matching that realizes the Fréchet distance, the two subcurves $P_{null}$ get roughly matched to different subcurves $Q_{null}$. Then, the subcurve in between the two points matched to $P_{null}$ contains $(Q_{1})^{d+1}$ as a subcurve. Now, we enforce that $(Q_{1})^{d+1}$ gets roughly matched to $(P_{0})^{d+1}$. Hence, afterwards $\bigcirc_{k=1}^{d}P_{v_{ik}}$ and $\bigcirc_{\ell=1}^{d}Q_{u_{j\ell}}$ have to be matched together for some $i$ and $j$. We will show that this is only possible if $\langle v_{i},u_{j}\rangle=0$. This matching is schematically visualized in Figure 1.
\par
\par
% END S3.p2
% BEGIN S3.F1
\begin{center}\includegraphics[width=0.7\linewidth,height=0.38\textheight,keepaspectratio]{LowerBoundSchematicView-page-2.pdf2svg.pdf}\end{center}
 Figure 1: A schematic view of the construction of the curves $P$ and $Q$. Here, $Q_{null}$ is drawn in red, $P_{0}$ and $Q_{0}$ in green, $P_{1}$ and $Q_{1}$ in blue, and $Q_{c}$ in orange. 
% END S3.F1
% BEGIN S3.p3
We define
\par
\[\begin{aligned}
\displaystyle P_{s}=\langle(-1,0)\rangle, & \displaystyle P_{null}=\langle(-5,0)\rangle, & \displaystyle P^{*}=\langle(-3,0)\rangle\circ\langle P_{0}\rangle^{2d+3}\circ%
\langle(-3,0)\rangle, \\
\displaystyle Q_{c}=\langle(0,0)\rangle, & \displaystyle Q_{null}=\langle(-4,0)\rangle, & \displaystyle Q^{*}=\langle(-2,0),(0,0),(-2,0)\rangle,
\end{aligned}\]
for 2D and the definition carries over directly to 1D. The curves $P_{0}$, $P_{1}$, $Q_{0}$, and $Q_{1}$ are visualized in Figure 2. For the 1D setting, let
\par
\[\begin{aligned}
\displaystyle P_{0}=\langle-1,1\rangle, & \displaystyle P_{1}=\langle-1,0\rangle, & \displaystyle Q_{0}=\langle-2,1\rangle, & \displaystyle Q_{1}=\langle-2,2\rangle.
\end{aligned}\]
\par
% END S3.p3
% BEGIN S3.p4
For the 2D constructions, define eight points with distance $1$ to $(0,0)$ where $a_{1}=(-1,0)$ and the direction of $a_{i+1}$ is the direction of $a_{i}$ rotated by $\pi/4$, i.e., in $L_{2}$
\par
\[\begin{aligned}
\displaystyle a_{1}=(-1,0), & \displaystyle a_{2}=(-1/\sqrt{2},-1/\sqrt{2}), & \displaystyle a_{3}=(0,-1), & \displaystyle a_{4}=(1/\sqrt{2},-1/\sqrt{2}), \\
\displaystyle a_{5}=(1,0), & \displaystyle a_{6}=(1/\sqrt{2},1/\sqrt{2}), & \displaystyle a_{7}=(0,1), & \displaystyle a_{8}=(-1/\sqrt{2},1/\sqrt{2}),
\end{aligned}\]
and in $L_{\infty}$
\par
\[\begin{aligned}
\displaystyle a_{1}=(-1,0), & \displaystyle a_{2}=(-1,-1), & \displaystyle a_{3}=(0,-1), & \displaystyle a_{4}=(1,-1), \\
\displaystyle a_{5}=(1,0), & \displaystyle a_{6}=(1,1), & \displaystyle a_{7}=(0,1), & \displaystyle a_{8}=(-1,1).
\end{aligned}\]
Define $c=\sqrt{2}$ and $\tilde{a}=a_{8}+a_{7}$ in $L_{2}$ and $c=2$ and $\tilde{a}=2a_{8}$ in $L_{\infty}$. Note that $c$ is defined such that $ca_{1}$ is the point with distance 1 to $a_{8}$ and $a_{2}$ that maximizes the distance to $(0,0)$ and $\tilde{a}$ such that $\|\tilde{a}-a_{2}\|=1+c$ and $\|\tilde{a}-a_{8}\|=1$. Then we set
\par
\[\begin{aligned}
\displaystyle P_{0}=\langle a_{1},a_{2},\dots a_{7},a_{8},a_{7}\dots,a_{1}\rangle, & \displaystyle P_{1}=\langle a_{1},a_{2},\dots,a_{8},a_{1},a_{2},a_{1}a_{8},%
\dots,a_{1}\rangle, \\
\displaystyle Q_{0}=\langle 2a_{1},2a_{2},\dots,2a_{8},ca_{1},2a_{8},\dots,2a_%
{1}\rangle, & \displaystyle Q_{1}=\langle 2a_{1},2a_{2},\dots,2a_{7},\tilde{a},2a_{7},\dots,%
2a_{1}\rangle.
\end{aligned}\]
We show that this results in an inapproximability result of $(1+c-\varepsilon)$ for time in $\mathcal{O}((nm)^{1-\delta})$.
\par
\par
% END S3.p4
% BEGIN S3.F2
\begin{center}\includegraphics[width=0.7\linewidth,height=0.38\textheight,keepaspectratio]{Lower_Bound_Curves-page-10.pdf2svg.pdf}\end{center}
\begin{center}\includegraphics[width=0.7\linewidth,height=0.38\textheight,keepaspectratio]{Lower_Bound_Curves-page-11.pdf2svg.pdf}\end{center}
\begin{center}\includegraphics[width=0.7\linewidth,height=0.38\textheight,keepaspectratio]{Lower_Bound_Curves-page-12.pdf2svg.pdf}\end{center}
\begin{center}\includegraphics[width=0.7\linewidth,height=0.38\textheight,keepaspectratio]{Lower_Bound_Curves-page-13.pdf2svg.pdf}\end{center}
\begin{center}\includegraphics[width=0.7\linewidth,height=0.38\textheight,keepaspectratio]{Lower_Bound_Curves-page-1.pdf2svg.pdf}\end{center}
\begin{center}\includegraphics[width=0.7\linewidth,height=0.38\textheight,keepaspectratio]{Lower_Bound_Curves-page-2.pdf2svg.pdf}\end{center}
\begin{center}\includegraphics[width=0.7\linewidth,height=0.38\textheight,keepaspectratio]{Lower_Bound_Curves-page-3.pdf2svg.pdf}\end{center}
\begin{center}\includegraphics[width=0.7\linewidth,height=0.38\textheight,keepaspectratio]{Lower_Bound_Curves-page-4.pdf2svg.pdf}\end{center}
\begin{center}\includegraphics[width=0.7\linewidth,height=0.38\textheight,keepaspectratio]{Lower_Bound_Curves-page-5.pdf2svg.pdf}\end{center}
\begin{center}\includegraphics[width=0.7\linewidth,height=0.38\textheight,keepaspectratio]{Lower_Bound_Curves-page-6.pdf2svg.pdf}\end{center}
\begin{center}\includegraphics[width=0.7\linewidth,height=0.38\textheight,keepaspectratio]{Lower_Bound_Curves-page-7.pdf2svg.pdf}\end{center}
\begin{center}\includegraphics[width=0.7\linewidth,height=0.38\textheight,keepaspectratio]{Lower_Bound_Curves-page-8.pdf2svg.pdf}\end{center}
  Figure 2: The curves $P_{0}$, $P_{1}$, $Q_{0}$, and $Q_{1}$. The upper curves are defined in 1D, the middle curves for the $L_{2}$-metric in 2D and the lower curves for the $L_{\infty}$-metric in 2D. 
% END S3.F2
% BEGIN Thmtheorem4
\par\medskip\noindent\textbf{Lemma 4.}\par
If there exist indices $i$ and $j$ such that $\langle v_{i},u_{j}\rangle=0$, then in every construction for $P$ and $Q$ above it holds that $d_{\text{dF}}(P,Q)\leq 1$.
\par
\par
\par
% END Thmtheorem4
% BEGIN S3.SS0.SSS0.P0.SPx1
\par\medskip\noindent\textbf{Proof.}\par
% BEGIN S3.SS0.SSS0.P0.SPx1.p1
For every construction of $P$ and $Q$, it holds that $d_{\text{dF}}(P^{*},Q^{*})\leq 1$, $d_{\text{dF}}(P_{s},Q^{*})\leq 1$. Further, for $x,y\in\{0,1\}$ with $\langle x,y\rangle=0$, it holds that $d_{\text{dF}}(P_{x},Q_{y})\leq 1$. Then, we can match the following with Fréchet distance at most 1 and the lemma follows by Observation 2
\par
\par
% END S3.SS0.SSS0.P0.SPx1.p1
% BEGIN S3.SS0.SSS0.P0.SPx1.p2
\[\begin{aligned}
\blacksquare & \displaystyle P_{s}\circ(P^{*})^{m-j} & \displaystyle(Q^{*})^{m-1}, \\
\blacksquare & \displaystyle(P^{*})^{j-1}\circ P_{null} & \displaystyle\bigcirc_{k=1}^{j-1}(Q_{null}\circ Q_{c}\circ(Q_{1})^{d+1}\circ Q%
_{0}\circ \\
\displaystyle\nobreak\ \nobreak\ \nobreak\ \nobreak\ \nobreak\ \nobreak\ %
\nobreak\ \nobreak\ \bigcirc_{\ell=1}^{d}Q_{u_{k\ell}}\circ Q_{0}\circ Q_{c}){%
\circ Q_{null}},\nobreak\ \nobreak\ \nobreak\ \nobreak\ \nobreak\  \\
\blacksquare & \displaystyle\bigcirc_{k=1}^{i-1}\left((P_{0})^{d+1}\circ P_{1}\circ\bigcirc_{%
\ell=1}^{d}P_{v_{k\ell}}\circ P_{1}\right) & \displaystyle Q_{c}, \\
\blacksquare & \displaystyle(P_{0})^{d+1}\circ P_{1}\circ\bigcirc_{\ell=1}^{d}P_{v_{i\ell}}%
\circ P_{1} & \displaystyle(Q_{1})^{d+1}\circ Q_{0}\circ\bigcirc_{\ell=1}^{d}Q_{u_{j\ell}}%
\circ Q_{0}, \\
\blacksquare & \displaystyle\bigcirc_{k=i+1}^{n}\left((P_{0})^{d+1}\circ P_{1}\circ\bigcirc_{%
\ell=1}^{d}P_{v_{k\ell}}\circ P_{1}\right) & \displaystyle Q_{c}, \\
\blacksquare & \displaystyle P_{null}\circ(P^{*})^{m-j} & \displaystyle\bigcirc_{k=j+1}^{m}(Q_{null}\circ Q_{c}\circ(Q_{1})^{d+1}\circ Q%
_{0}\circ \\
\displaystyle\nobreak\ \nobreak\ \nobreak\ \nobreak\ \nobreak\ \nobreak\ %
\nobreak\ \nobreak\ \bigcirc_{\ell=1}^{d}Q_{u_{k\ell}}\circ Q_{0}\circ Q_{c}){%
\circ Q_{null}}, \\
\blacksquare & \displaystyle(P^{*})^{j-1}\circ P_{s} & \displaystyle(Q^{*})^{m-1}.\ 
\end{aligned}\]
$\hfill\blacktriangleleft$
\par
\par
% END S3.SS0.SSS0.P0.SPx1.p2
% BEGIN S3.SS0.SSS0.P0.SPx1.p3
To show the other direction, we first prove some lemmas.
\par
\par
% END S3.SS0.SSS0.P0.SPx1.p3
% BEGIN Thmtheorem5
\par\medskip\noindent\textbf{Lemma 5.}\par
Consider the 1D setting and let $Q^{\prime}=\bigcirc_{\ell=1}^{d^{\prime}}Q_{\sigma_{\ell}}$, where $\sigma_{\ell}\in\{0,1\}$ for every $\ell$. Further, let $s\leq t$ be such that $P[s,t]$ is a subcurve in between the two subcurves $P_{null}$ and $d_{\text{dF}}(Q^{\prime},P[s,t])<2$. Then $P[s,t]=\bigcirc_{\ell=1}^{d^{\prime}}P_{\pi_{\ell}}$ where $\pi_{\ell}=0$ if $\sigma_{\ell}=1$ and $\pi_{\ell}\in\{0,1\}$ if $\sigma_{\ell}=0$ for every $\ell$.
\par
\par
\par
% END Thmtheorem5
% END S3.SS0.SSS0.P0.SPx1
% BEGIN S3.SS0.SSS0.P0.SPx2
\par\medskip\noindent\textbf{Proof.}\par
% BEGIN S3.SS0.SSS0.P0.SPx2.p1
It holds that every second vertex of $Q^{\prime}$ is $-2$ and the vertices of value $-1$ on $P[s,t]$ are the only vertices on $P[s,t]$ within distance less than $2$ to those vertices of $Q^{\prime}$. Further, it holds that the vertices with value $-1$ of $P[s,t]$ have distance less than two only to the vertices of value $-2$ of $Q^{\prime}$. Hence, since the vertices of value $-1$ on $P[s,t]$ are exactly every second vertex on $P[s,t]$, we must traverse $P[s,t]$ and $Q^{\prime}$ synchronously. Therefore, if $\sigma_{\ell}=0$ then $\pi_{\ell}\in\{0,1\}$ and if $\sigma_{\ell}=1$ then $\pi_{\ell}=0$ by construction. $\hfill\blacktriangleleft$
\par
\par
% END S3.SS0.SSS0.P0.SPx2.p1
% BEGIN S3.SS0.SSS0.P0.SPx2.p2
In 2D, we also want to show a bound for the continuous Fréchet distance. Then the vertices of $Q$ can also be matched to points on edges of $P$. We use the notation of the following subcurve of $P_{0}$ (resp. $P_{1}$): $P_{0}^{\prime}=\langle a_{4},a_{5},\dots,a_{8},\dots,a_{4}\rangle$ (resp. $P_{1}^{\prime}=\langle a_{4},\dots,a_{8},a_{1},a_{2},a_{1},a_{8},\dots,a_{4}\rangle$).
\par
\par
% END S3.SS0.SSS0.P0.SPx2.p2
% BEGIN Thmtheorem6
\par\medskip\noindent\textbf{Lemma 6.}\par
Consider the 2D setting. For any subcurve $P^{\prime}$ of $P$ in between the two subcurves $P_{null}$ with $d_{\text{F}}(P^{\prime},Q_{y})<1+c$ for $y\in\{0,1\}$, it holds that $P^{\prime}_{x}$ is a subcurve of $P^{\prime}$ and $P^{\prime}$ is a subcurve of $\langle a_{4},a_{3},a_{2}\rangle\circ P_{x}\circ\langle a_{2},a_{3},a_{4}\rangle$, where $x=0$ if $y=1$ and $x\in\{0,1\}$ if $y=0$.
\par
\par
\par
% END Thmtheorem6
% BEGIN S3.F3
\begin{center}\includegraphics[width=0.7\linewidth,height=0.38\textheight,keepaspectratio]{Lower_Bound_Curves-page-9.pdf2svg.pdf}\end{center}
 Figure 3: In white: the non-free space of the curve $Q_{0}$ and $P_{0}\circ P_{0}\circ P_{1}$ on the left and of the curve $Q_{1}$ and $P_{0}\circ P_{0}\circ P_{1}$ on the right. The red, orange and pink paths visualize the proof of Lemma 6. 
% END S3.F3
% END S3.SS0.SSS0.P0.SPx2
% BEGIN S3.SS0.SSS0.P0.SPx3
\par\medskip\noindent\textbf{Proof.}\par
% BEGIN S3.SS0.SSS0.P0.SPx3.p1
We show this lemma by constructing paths in the non-free space for $\delta=1+c$ (see Figure 3). First observe that $\|2a_{i}-a_{j}\|\geq 1+c$ for any $i,j\in\{1,\dots,8\}$ with ${i-j\geq 3\mod 8}$.
\par
\par
% END S3.SS0.SSS0.P0.SPx3.p1
% BEGIN S3.SS0.SSS0.P0.SPx3.p2
Assume that the first (resp. last) vertex of $Q_{y}$ is matched to a point on a subcurve $P^{\prime}_{x}$ with $x\in\{0,1\}$. Then, there must exist a path completely contained in the free space ending at the last (resp. first) vertex of $Q_{y}$. This cannot be the case since there exists a continuous matching between $P^{\prime}_{x}$ and $\langle 2a_{1},2a_{2}\dots,2a_{5},\dots,2a_{1}\rangle$ that is completely contained in the non-free space (see the red paths in Figure 3). Therefore, the first as well as the last vertex of $Q_{y}$ has to be matched to a point on a subcurve $P^{\prime\prime}$ of $P$ such that $P^{\prime\prime}=\langle a_{4},a_{3},a_{2},a_{1},a_{2},a_{3},a_{4}\rangle$.
\par
\par
% END S3.SS0.SSS0.P0.SPx3.p2
% BEGIN S3.SS0.SSS0.P0.SPx3.p3
Next, we show that $P^{\prime}$ contains exactly one $P^{\prime}_{x}$ as a subcurve. Since $\|ca_{1}-a_{5}\|\geq 1+c$ (resp. $\|\tilde{a}-a_{5}\|\geq 1+c$), there exists a path that is completely contained in the non-free space between $\langle a_{5},a_{4},\dots a_{2},a_{1},a_{2},\dots,a_{4},a_{5}\rangle$ and $\langle 2a_{8},ca_{1},2a_{8},2a_{7},\dots,2a_{4},2a_{5},\dots,2a_{8},ca_{1},2a%
_{8}\rangle$(resp. $\langle\tilde{a},2a_{7},\dots,2a_{4},2a_{5},\dots,\tilde{a}\rangle$) (see the orange paths in Figure 3). Hence by construction of $P$, it follows that $P^{\prime}_{x}$ is a subcurve of $P^{\prime}$ and $P^{\prime}$ is a subcurve of $\langle a_{4},a_{3},a_{2}\rangle\circ P_{x}\circ\langle a_{2},a_{3},a_{4}\rangle$ with $x\in\{0,1\}$. It remains to show that $x\neq 1$, if $y=1$. This follows since the distance between $a_{2}$ and $2a_{7}$, $\tilde{a}$, $2a_{7}$ is at least $1+c$ (see the pink path in Figure 3). $\hfill\blacktriangleleft$ The next lemma follows directly by the construction of the curves and Lemma 6.
\par
\par
% END S3.SS0.SSS0.P0.SPx3.p3
% BEGIN Thmtheorem7
\par\medskip\noindent\textbf{Lemma 7.}\par
Consider the 2D setting and let $Q^{\prime}=\bigcirc_{\ell=1}^{d^{\prime}}Q_{\sigma_{\ell}}$, where $\sigma_{\ell}\in\{0,1\}$ for every $\ell$. Further, let $s\leq t$ be such that $P[s,t]$ is a subcurve in between the two subcurves $P_{null}$ and $d_{\text{F}}(Q^{\prime},P[s,t])<1+c$. Then, $P_{\pi_{1}}^{\prime}\circ\bigcirc_{\ell=2}^{d^{\prime}-1}P_{\pi_{\ell}}\circ P%
_{\pi_{d^{\prime}}}^{\prime}$ is a subcurve of $P[s,t]$ and $P[s,t]$ is a subcurve of $\langle a_{4},a_{3},a_{2}\rangle\circ\bigcirc_{\ell=1}^{d^{\prime}}P_{\pi_{%
\ell}}\circ\langle a_{2},a_{3},a_{4}\rangle$ where $\pi_{\ell}=0$ if $\sigma_{\ell}=1$ and $\pi_{\ell}\in\{0,1\}$ if $\sigma_{\ell}=0$ for every $\ell$.
\par
\par
\par
% END Thmtheorem7
% BEGIN Thmtheorem8
\par\medskip\noindent\textbf{Lemma 8.}\par
If $d_{\text{dF}}(P,Q)<2$ (resp. $d_{\text{F}}(P,Q)<1+c$) in the 1D (resp. 2D) construction, then there exist two indices $i$ and $j$ such that $\langle v_{i},u_{j}\rangle=0$.
\par
\par
\par
% END Thmtheorem8
% END S3.SS0.SSS0.P0.SPx3
% BEGIN S3.SS0.SSS0.P0.SPx4
\par\medskip\noindent\textbf{Proof.}\par
% BEGIN S3.SS0.SSS0.P0.SPx4.p1
Let $M$ be a discrete (resp. continuous) matching that realizes the discrete (resp. continuous) Fréchet distance between $P$ and $Q$. Assume that $d_{\text{F}}(P,Q)<2$ (resp. $d_{\text{dF}}(P,Q)<1+c$), then in the matching $M$, the two points of value $P_{null}$ of $P$ must be matched to two distinct points $q_{1},q_{2}$ on $Q$ both on an edge that contains the value $Q_{null}$, since they are the only points within distance at most $1+c\leq 3$ to $P_{null}$. Let $P_{0}^{\prime}=P_{0}$ and $P_{1}^{\prime}=P_{1}$ in the 1D setting. Define $j$ such that $q_{1}$ is contained on an incident edge to the $j$-th $Q_{null}$ on $Q$. Then, the first $(Q_{1})^{d+1}$ after the $j$-th $Q_{null}$ on $Q$ must be matched to a subcurve $P^{\prime}$ of $P$ such that $P_{1}^{\prime}\circ(P_{0})^{d+1}\circ P_{1}^{\prime}$ is a subcurve of $P^{\prime}$ by Lemma 5 and Lemma 7. The only time we have $(P_{0})^{d+1}$ in $P$ is before $P_{1}\circ\bigcirc_{\ell=1}^{d}P_{v_{i\ell}}$ for an $i$. Hence by Lemma 5 and Lemma 7, there exists an index $i$ such that in the matching $M$ the subcurve $P^{\prime}_{v_{i,1}}\circ\bigcirc_{\ell=2}^{d-1}P_{v_{i,\ell}}\circ P^{\prime}%
_{v_{i,d}}$ is matched to a subcurve of $\bigcirc_{\ell=1}^{d}Q_{u_{j,\ell}}$ and $v_{i,\ell}=0$ if ${u_{j,\ell}}=1$. Hence, it holds that $\langle v_{i},u_{j}\rangle=0$. $\hfill\blacktriangleleft$ By Lemma 4 and Lemma 8 together with Observation 3, the next two theorems follow.
\par
\par
% END S3.SS0.SSS0.P0.SPx4.p1
% BEGIN Thmtheorem9
\par\medskip\noindent\textbf{Theorem 9.}\par
For any $\delta,\varepsilon>0$, there does not exist a $(2-\varepsilon)$-approximation algorithm for the discrete Fréchet distance between 1-dimensional curves of complexity $n$ and $m$ using $\mathcal{O}((nm)^{1-\delta})$ time unless OVH fails.
\par
\par
\par
% END Thmtheorem9
% BEGIN Thmtheorem10
\par\medskip\noindent\textbf{Theorem 10.}\par
For any $\delta,\varepsilon>0$, there does not exist a $(k-\varepsilon)$-approximation algorithm for the continuous or discrete Fréchet distance between 2-dimensional curves of complexity $n$ and $m$ using $\mathcal{O}((nm)^{1-\delta})$ time unless OVH fails, where $k=3$ (resp. $k=1+\sqrt{2}\approx 2.41$) if the curves lie in the $L_{\infty}$-space (resp. in the Euclidean space).
\par
\par
\par
% END Thmtheorem10
% BEGIN S3.SS0.SSS0.P0.SPx4.p2
Our construction gives also a bound on Fréchet distance oracles (proof in the full version).
\par
\par
% END S3.SS0.SSS0.P0.SPx4.p2
% END S3.SS0.SSS0.P0.SPx4
% END S3
% BEGIN S5
\par\medskip\noindent\textbf{5 $(3+\varepsilon)$-approximation Algorithm for the Fréchet distance}\par
% BEGIN S5.p1
In this section, we give a $(3+\varepsilon)$-approximation algorithm to compute the continuous/discrete Fréchet distance between two curves $P$ and $Q$ of complexity $n$ resp. $m$ with $m\leq n$. We describe this section for the continuous Fréchet distance, but everything works analogously for the discrete Fréchet distance as well. We begin with the decision variant for $\delta$. In the first step, using the curve $Q$, we compute a curve $P^{\prime}$ of complexity at most $2m$ such that $d_{\text{F}}(P,P^{\prime})\leq\delta$ or decide that $d_{\text{F}}(P,Q)>\delta$ in linear time.
\par
\par
% END S5.p1
% BEGIN Thmtheorem21
\par\medskip\noindent\textbf{Lemma 21.}\par
Algorithm 3 computes a curve $P^{\prime}$ of complexity at most $2m$ with $d_{\text{F}}(P,P^{\prime})\leq\delta$ or returns $d_{\text{F}}(P,Q)>\delta$ in $\mathcal{O}(n)$ time.
\par
\par
\par
% END Thmtheorem21
% BEGIN S5.SS0.SSS0.P0.SPx1
\par\medskip\noindent\textbf{Proof.}\par
% BEGIN S5.SS0.SSS0.P0.SPx1.p1
In the $i$-th iteration of the while-loop, it holds that $s_{i}\geq i$ by line 3 of Algorithm 3. Further, it holds that $a_{i}\leq a_{i+1}$ for every $i$ and line 6 has running time $\mathcal{O}(a_{i+1}-a_{i})$ as $\overline{Q(s_{i})Q(t_{i})}\subset\overline{Q(s_{i})Q(\lfloor s_{i}+1\rfloor)}$. Hence, the running time follows. If a curve $P^{\prime}$ is returned, it holds that $d_{\text{F}}(P,P^{\prime})\leq\delta$ by line 6 and Observation 2. The complexity of $P^{\prime}$ is at most $2m$ since there are at most $m$ iterations of the while loop. It remains to show correctness for the case when $d_{\text{F}}(P,Q)>\delta$ is returned. Assume that $d_{\text{F}}(P,Q)>\delta$ is returned and in fact it holds $d_{\text{F}}(P,Q)\leq\delta$. Let $M$ be the matching between $P$ and $Q$ that realizes the Fréchet distance. We consider two cases (see Figure 6). In the first case, there exist $b,b^{\prime}$ and an index $i$ such that $(b,s_{i}),(b^{\prime},t_{i})\in M$ with $b\leq a_{i}$ and $b^{\prime}>a_{i+1}$. This contradicts line 6 of the algorithm. Otherwise, there must exist an index $i$ and a parameter $c$ such that $(a_{i},c)\in M$ with $t_{i-1}\leq c<s_{i}$. This contradicts line 3. Hence, it follows that $d_{\text{F}}(P,Q)>\delta$. $\hfill\blacktriangleleft$
\par
\par
% END S5.SS0.SSS0.P0.SPx1.p1
% BEGIN S5.F6
\begin{center}\includegraphics[width=0.7\linewidth,height=0.38\textheight,keepaspectratio]{Simplification_depending_on_Q-page-1.pdf2svg.pdf}\end{center}
\begin{center}\includegraphics[width=0.7\linewidth,height=0.38\textheight,keepaspectratio]{Simplification_depending_on_Q-page-2.pdf2svg.pdf}\end{center}
  Figure 6: The green curve is a matching between the curves $P[a_{i},a_{i+1}]$ and $Q[s_{i},t_{i}]$ computed in Algorithm 3. The blue curves visualize the matching used in the proof of Lemma 21. 
% END S5.F6
% BEGIN algorithm3
Algorithm 3  Simplification of $P$ depending on $Q$. 
\begin{center}\includegraphics[width=0.7\linewidth,height=0.38\textheight,keepaspectratio]{dagpub-standalone-combined-2026-06-01-18-18-02-page-3.pdf2svg.pdf}\end{center}
% END algorithm3
% BEGIN S5.SS0.SSS0.P0.SPx1.p2
Lemma 21 can be easily adapted such that it holds for the discrete Fréchet distance. Given $P^{\prime}$, we decide whether $d_{\text{F}}(P^{\prime},Q)\leq 2\delta$ in $\mathcal{O}(m^{2})$ time [3] or slightly faster [15]. If $d_{\text{F}}(P^{\prime},Q)\leq 2\delta$, then $d_{\text{F}}(P,Q)\leq d_{\text{F}}(P^{\prime},Q)+d_{\text{F}}(P,P^{\prime})%
\leq 3\delta$. Otherwise, it holds that $d_{\text{F}}(P,Q)\geq d_{\text{F}}(P^{\prime},Q)-d_{\text{F}}(P,P^{\prime})>\delta$. Hence, we obtain the following lemma.
\par
\par
% END S5.SS0.SSS0.P0.SPx1.p2
% BEGIN Thmtheorem22
\par\medskip\noindent\textbf{Lemma 22.}\par
Given two curves $P$ and $Q$ of complexity $n$ and $m$ in any metric space $M$. Assume that we can compute distances in $M$ in constant time. We can $3$-approximate the decision variant of the continuous/discrete Fréchet distance in $\mathcal{O}(n+m^{2})$ time.
\par
\par
\par
% END Thmtheorem22
% BEGIN S5.SS0.SSS0.P0.SPx1.p3
For curves in the Euclidean space, Colombe and Fox (Theorem 11 of [18]) show how to transform a decision algorithm into an approximate optimization algorithm increasing the running time by a $\mathcal{O}(\log(n/\varepsilon))$ factor. Since all $L_{p}$-norms on $\mathbb{R}^{d}$ are equivalent and Lemma 22 holds for all $L_{p}$-norms, we can adopt their approach such that we can use it for all $L_{p}$-norms.
\par
\par
% END S5.SS0.SSS0.P0.SPx1.p3
% BEGIN Thmtheorem23
\par\medskip\noindent\textbf{Theorem 23.}\par
There exist a $(3+\varepsilon)$-approximation algorithm for the continuous/discrete Fréchet distance using $\mathcal{O}((n+m^{2})\log(n/\varepsilon))$ time.
\par
\par
\par
% END Thmtheorem23
% END S5.SS0.SSS0.P0.SPx1
% END S5
% BEGIN bib
\par\medskip\noindent\textbf{References}\par
\begin{enumerate}
\item[]  [1]  Pankaj K. Agarwal, Rinat Ben Avraham, Haim Kaplan, and Micha Sharir.  Computing the discrete Fréchet distance in subquadratic time.  SIAM Journal on Computing, 43(2):429–449, 2014.  doi:10.1137/130920526.  
\item[]  [2]  Pankaj K. Agarwal, Sariel Har-Peled, Nabil H. Mustafa, and Yusu Wang.  Near-linear time approximation algorithms for curve simplification.  Algorithmica, 42(3-4):203–219, 2005.  doi:10.1007/S00453-005-1165-Y.  
\item[]  [3]  Helmut Alt and Michael Godau.  Computing the Fréchet distance between two polygonal curves.  Int. J. Comput. Geom. Appl., 5:75–91, 1995.  doi:10.1142/S0218195995000064.  
\item[]  [4]  Sergey Bereg, Minghui Jiang, Wencheng Wang, Boting Yang, and Binhai Zhu.  Simplifying 3d polygonal chains under the discrete Fréchet distance.  In Eduardo Sany Laber, Claudson F. Bornstein, Loana Tito Nogueira, and Luérbio Faria, editors, LATIN 2008: Theoretical Informatics, 8th Latin American Symposium, Búzios, Brazil, April 7-11, 2008, Proceedings, volume 4957 of Lecture Notes in Computer Science, pages 630–641. Springer, 2008.  doi:10.1007/978-3-540-78773-0\_54.  
\item[]  [5]  Lotte Blank and Anne Driemel.  A faster algorithm for the Fréchet distance in 1d for the imbalanced case.  In Timothy M. Chan, Johannes Fischer, John Iacono, and Grzegorz Herman, editors, 32nd Annual European Symposium on Algorithms, ESA 2024, September 2-4, 2024, Royal Holloway, London, United Kingdom, volume 308 of LIPIcs, pages 28:1–28:15. Schloss Dagstuhl – Leibniz-Zentrum für Informatik, 2024.  doi:10.4230/LIPIcs.ESA.2024.28.  
\item[]  [6]  Karl Bringmann.  Why walking the dog takes time: Fréchet distance has no strongly subquadratic algorithms unless seth fails.  In 2014 IEEE 55th Annual Symposium on Foundations of Computer Science, pages 661–670. IEEE, 2014.  doi:10.1109/FOCS.2014.76.  
\item[]  [7]  Karl Bringmann and Bhaskar Ray Chaudhury.  Polyline simplification has cubic complexity.  In Gill Barequet and Yusu Wang, editors, 35th International Symposium on Computational Geometry, SoCG 2019, June 18-21, 2019, Portland, Oregon, USA, volume 129 of LIPIcs, pages 18:1–18:16. Schloss Dagstuhl – Leibniz-Zentrum für Informatik, 2019.  doi:10.4230/LIPIcs.SOCG.2019.18.  
\item[]  [8]  Karl Bringmann, Anne Driemel, André Nusser, and Ioannis Psarros.  Tight bounds for approximate near neighbor searching for time series under the Fréchet distance.  In Joseph (Seffi) Naor and Niv Buchbinder, editors, Proceedings of the 2022 ACM-SIAM Symposium on Discrete Algorithms, SODA 2022, Virtual Conference / Alexandria, VA, USA, January 9 - 12, 2022, pages 517–550. SIAM, 2022.  doi:10.1137/1.9781611977073.25.  
\item[]  [9]  Karl Bringmann and Marvin Künnemann.  Multivariate fine-grained complexity of longest common subsequence.  In Artur Czumaj, editor, Proceedings of the Twenty-Ninth Annual ACM-SIAM Symposium on Discrete Algorithms, SODA 2018, New Orleans, LA, USA, January 7-10, 2018, pages 1216–1235. SIAM, 2018.  doi:10.1137/1.9781611975031.79.  
\item[]  [10]  Karl Bringmann and Wolfgang Mulzer.  Approximability of the discrete Fréchet distance.  Journal of Computational Geometry, 7(2):46–76, 2016.  doi:10.20382/JOCG.V7I2A4.  
\item[]  [11]  Kevin Buchin, Maike Buchin, Wouter Meulemans, and Wolfgang Mulzer.  Four soviets walk the dog: Improved bounds for computing the Fréchet distance.  Discrete Comput. Geom., 58(1):180–216, July 2017.  doi:10.1007/s00454-017-9878-7.  
\item[]  [12]  Kevin Buchin, Tim Ophelders, and Bettina Speckmann.  SETH says: Weak Fréchet distance is faster, but only if it is continuous and in one dimension.  In Timothy M. Chan, editor, Proceedings of the Thirtieth Annual ACM-SIAM Symposium on Discrete Algorithms, SODA 2019, San Diego, California, USA, January 6-9, 2019, pages 2887–2901. SIAM, 2019.  doi:10.1137/1.9781611975482.179.  
\item[]  [13]  Siu-Wing Cheng and Haoqiang Huang.  Curve simplification and clustering under Fréchet distance.  In Nikhil Bansal and Viswanath Nagarajan, editors, Proceedings of the 2023 ACM-SIAM Symposium on Discrete Algorithms, SODA 2023, Florence, Italy, January 22-25, 2023, pages 1414–1432. SIAM, 2023.  doi:10.1137/1.9781611977554.CH51.  
\item[]  [14]  Siu-Wing Cheng and Haoqiang Huang.  Solving Fréchet distance problems by algebraic geometric methods.  In David P. Woodruff, editor, Proceedings of the 2024 ACM-SIAM Symposium on Discrete Algorithms, SODA 2024, Alexandria, VA, USA, January 7-10, 2024, pages 4502–4513. SIAM, 2024.  doi:10.1137/1.9781611977912.158.  
\item[]  [15]  Siu-Wing Cheng and Haoqiang Huang.  Fréchet distance in subquadratic time.  In Yossi Azar and Debmalya Panigrahi, editors, Proceedings of the 2025 Annual ACM-SIAM Symposium on Discrete Algorithms, SODA 2025, New Orleans, LA, USA, January 12-15, 2025, pages 5100–5113. SIAM, 2025.  doi:10.1137/1.9781611978322.173.  
\item[]  [16]  Siu-Wing Cheng, Haoqiang Huang, and Le Jiang.  Simplification of trajectory streams.  In Oswin Aichholzer and Haitao Wang, editors, 41st International Symposium on Computational Geometry, SoCG 2025, June 23-27, 2025, Kanazawa, Japan, volume 332 of LIPIcs, pages 34:1–34:14. Schloss Dagstuhl – Leibniz-Zentrum für Informatik, 2025.  doi:10.4230/LIPIcs.SOCG.2025.34.  
\item[]  [17]  Siu-Wing Cheng, Haoqiang Huang, and Shuo Zhang.  Constant approximation of Fréchet distance in strongly subquadratic time.  In Michal Koucký and Nikhil Bansal, editors, Proceedings of the 57th Annual ACM Symposium on Theory of Computing, STOC 2025, Prague, Czechia, June 23-27, 2025, pages 2329–2340. ACM, 2025.  doi:10.1145/3717823.3718157.  
\item[]  [18]  Connor Colombe and Kyle Fox.  Approximating the (continuous) Fréchet distance.  In Kevin Buchin and Éric Colin de Verdière, editors, 37th International Symposium on Computational Geometry, SoCG 2021, June 7-11, 2021, Buffalo, NY, USA (Virtual Conference), volume 189 of LIPIcs, pages 26:1–26:14. Schloss Dagstuhl – Leibniz-Zentrum für Informatik, 2021.  doi:10.4230/LIPIcs.SOCG.2021.26.  
\item[]  [19]  Anne Driemel and Sariel Har-Peled.  Jaywalking your dog: Computing the Fréchet distance with shortcuts.  SIAM J. Comput., 42(5):1830–1866, 2013.  doi:10.1137/120865112.  
\item[]  [20]  Anne Driemel, Ioannis Psarros, and Melanie Schmidt.  Sublinear data structures for short Fréchet queries.  CoRR, abs/1907.04420, 2019.  arXiv:1907.04420.  
\item[]  [21]  Thomas Eiter and Heikki Mannila.  Computing discrete Fréchet distance, 1994.  
\item[]  [22]  Arnold Filtser and Omrit Filtser.  Static and streaming data structures for Fréchet distance queries.  ACM Trans. Algorithms, 19(4):39:1–39:36, 2023.  doi:10.1145/3610227.  
\item[]  [23]  Greg N. Frederickson and Donald B. Johnson.  Generalized selection and ranking: Sorted matrices.  SIAM J. Comput., 13(1):14–30, 1984.  doi:10.1137/0213002.  
\item[]  [24]  Joachim Gudmundsson, Martin P. Seybold, and Sampson Wong.  Map matching queries on realistic input graphs under the Fréchet distance.  ACM Trans. Algorithms, 20(2):14, 2024.  doi:10.1145/3643683.  
\item[]  [25]  Aviad Rubinstein.  Hardness of approximate nearest neighbor search.  In Ilias Diakonikolas, David Kempe, and Monika Henzinger, editors, Proceedings of the 50th Annual ACM SIGACT Symposium on Theory of Computing, STOC 2018, Los Angeles, CA, USA, June 25-29, 2018, pages 1260–1268. ACM, 2018.  doi:10.1145/3188745.3188916.  
\item[]  [26]  Thijs van der Horst and Tim Ophelders.  Faster Fréchet distance approximation through truncated smoothing.  CoRR, abs/2401.14815, 2024.  doi:10.48550/arXiv.2401.14815.  
\item[]  [27]  Thijs van der Horst, Marc J. van Kreveld, Tim Ophelders, and Bettina Speckmann.  A subquadratic n${}^{\mbox{{$\epsilon$}}}$-approximation for the continuous Fréchet distance.  In Nikhil Bansal and Viswanath Nagarajan, editors, Proceedings of the 2023 ACM-SIAM Symposium on Discrete Algorithms, SODA 2023, Florence, Italy, January 22-25, 2023, pages 1759–1776. SIAM, 2023.  doi:10.1137/1.9781611977554.CH67.  
\item[]  [28]  Marc J. van Kreveld, Maarten Löffler, and Lionov Wiratma.  On optimal polyline simplification using the Hausdorff and Fréchet distance.  In Bettina Speckmann and Csaba D. Tóth, editors, 34th International Symposium on Computational Geometry, SoCG 2018, June 11-14, 2018, Budapest, Hungary, volume 99 of LIPIcs, pages 56:1–56:14. Schloss Dagstuhl – Leibniz-Zentrum für Informatik, 2018.  doi:10.4230/LIPIcs.SOCG.2018.56.  
\item[]  [29]  Ryan Williams.  A new algorithm for optimal 2-constraint satisfaction and its implications.  Theor. Comput. Sci., 348(2-3):357–365, 2005.  doi:10.1016/J.TCS.2005.09.023.  
\end{enumerate}
% END bib
